\scanpage{19}
\mspage{18}

De plus, pour une section $g$ de $X/S$, on vérifie que l'on a pour $\beta\deleted{\unsure{\in A}}\in\underline{\mathrm{Pic}}^0_{X/S}(S)=B(S)$

\[
l_g\bigl(\underset{\underset{\underline{\mathrm{Pic}}^\delta_{X/S}(S)}{\rotatebox{90}{$\in$}}}{\lambda}+\beta\bigr)=l_g(\lambda)\,\deleted{\unsure{+}}\unsure{-|+}\,\widetilde{\varphi(\delta)}(\beta)\tag*{**}
\]

\deleted{En effet, on sait que (comme tt \unsure{$S$-hom.|flèche} de} \deleted{torseurs sous des schémas abéliens différents)}

\[
\deleted{l_g^\delta\colon\underset{\text{torseur sous }B}{\underline{\mathrm{Pic}}^\delta_{X/S}}\longrightarrow\underset{\text{torseur sous }A}{\underline{\mathrm{Pic}}^{\varphi(\delta)}_{B/S}}\qquad\text{est associé}}
\]

\deleted{à un hom $\rho\colon B\longrightarrow A$, et il faut prouver que} \deleted{$\rho=-\widetilde{\varphi(\delta)}$. On peut pour ceci supposer $S$ le} \deleted{spectre d'un corps alg. clos\inserted{ $k$}, et il suffit de} \deleted{prouver que l'on a égalité pour les pts à} \deleted{valeurs dans $k$, i.e. \inserted{prouver} la formule ** pour $S=\operatorname{Spec}k$.} \deleted{Il suffit de la prouver pour des $\alpha$ engendrant} \deleted{le groupe $B(k)$, et donc pour $\alpha$ de la forme} \deleted{$\psi(g')-\psi(g)$ $(g',g\in X)$. \unsure{Choisissons un}}

On est ramené, grâce au fait que le faisceau inv. $L_\beta$ sur $X$ provient d'un faisceau \deleted{\unsure{unité associé} $L_\beta$}\inserted{$W_\beta$} sur $A$ (via $X\mathbin{\unsure{-|\to}}A$ déduit de $g$), sur \deleted{faisceau $W_\beta$ sur}\inserted{celui-ci défini par}
% slanted

\margin{\unsure{Vérif.|Réf.} ; Attention aux conventions de \unsure{signe|signes}}
le faisceau de Weil $W$ sur $A\times B$, au cas où $X=A$, $g=$ section nulle. Avec les notations développées ailleurs, on doit donc \unsure{identifier|expliciter|vérifier}

\[
\det\underline{\unsure{\Phi|\mathscr L|\mathfrak F}}_{\mathscr L}(L\otimes L_\beta)\simeq\det\Phi_{\mathscr L}(L)\otimes\underbrace{L'_{-\varphi(\delta)(\beta)}}_{\substack{\text{faisceau sur }B\text{ associé}\\\text{à }-\varphi(\delta)(\beta)}}\,,
\]
